At CERN, provisioning, configuration, and operational state are each handled by a different system. Since 2012 the IT department has worked to bridge them under the Agile Infrastructure project, combining OpenStack, Puppet, and Foreman with a set of in-house services so that users see one coherent platform.

The result is mature, but the seams between these systems are still where the work goes. Traceability and reproducibility suffer most when a service is deployed for the first time, scaled up, or rebuilt after a failure. Exposing newer cloud features, such as load balancers or security groups, means writing another integration each time.

Infrastructure-as-Code (IaC) tools such as Terraform and its open-source fork, OpenTofu, offer a declarative and version-controlled model that can sit across these procedures. They also bring an ecosystem of off-the-shelf modules that can be reused rather than rebuilt.

We present a framework that plugs into the existing infrastructure and configuration management landscape through declarative definitions, custom providers for in-house services, and CI/CD pipelines that validate every change before it is applied. Using a set of case studies, we show where the framework cuts operational complexity, how it helps in disaster recovery, and what it adds in automation, auditability, and replication across environments, offering a foundation for the evolution of infrastructure and configuration management at CERN.
